\textbf{Задание №3} Найти последовательность правил р, |p| > 10, которую можно применить при произвольном выводе терминальной цепочки, но нельзя применить при левом или правом выводе в заданной КС-грамматике (см. варианты заданий).
    Записать вывод v, в процессе которого применяется последовательность правил р. Построить дерево вывода.
    Записать левый и правый выводы, эквивалентные выводу v.\label{sec:3}
    
    \begin{enumerate}
    

    \item Произвольный вывод \(v\), в котором применяется \(p\)

    \[
    \begin{array}{ll@{\hspace{2em}}ll}
    1. & S \to Aba & 5. & A \to BaB \\
    2. & S \to bS  & 6. & A \to \varepsilon \\
    3. & S \to cAbA & 7. & B \to b \\
    4. & A \to AaS & 8. & B \to aA
    \end{array}
    \]

    \begin{align*}
    S_2 &\Rightarrow bS_2 \Rightarrow bbS_2 \Rightarrow bbbS_1 \Rightarrow bbbA_4ba 
        \Rightarrow bbbA_5aSba \\ 
        &\Rightarrow bbbBaBaS_2ba 
        \Rightarrow bbbBaB_7abSba 
        \Rightarrow bbbB_7ababSba \\
        &\Rightarrow bbbbababS_1ba 
        \Rightarrow bbbbababA_6baba 
        \Rightarrow bbbbababbaba
    \end{align*}

    Последовательность правил: \(2, 2, 2, 1, 4, 5, 2, 7, 7, 1, 6\).

    \item Дерево вывода
    
    \begin{center}
\begin{forest}
[$S_2$
  [b]
  [$S_2$
    [b]
    [$S_2$
      [b]
      [$S_1$
        [$A_4$
          [$A_5$
            [$B_7$ [b]]
            [a]
            [$B_7$ [b]]
          ]
          [a]
          [$S_2$
            [b]
            [$S_1$
              [$A_6$ [\(\varepsilon\)]]
              [b]
              [a]
            ]
          ]
        ]
        [b]
        [a]
      ]
    ]
  ]
]
\end{forest}
\end{center}
    
    \item Левый вывод
    \(2,2,2,1,4,5\)
    \begin{align*}
    S_2 &\Rightarrow bS_2 \Rightarrow bbS_2 \Rightarrow bbbS_1 \Rightarrow bbbA_4ba \Rightarrow bbbA_5aSba \Rightarrow bbbBaBaS_2ba 
    \end{align*}

    Крайний слева нетерминал --- это \(B\). Следующее правило в \(p\) ---
    \(2\) (\(S\to bS\)), но его нельзя применить к \(B\). Значит, левый вывод
    по этой последовательности невозможен.
    
    \item Правый вывод \(2,2,2,1,4\)
    \begin{align*}
    S_2 &\Rightarrow bS_2 \Rightarrow bbS_2 \Rightarrow bbbS_1 \Rightarrow bbbA_4ba \Rightarrow bbbA_5aSba 
    \end{align*}
    Крайний справа нетерминал --- это \(S\). Следующее правило в \(p\) ---
    \(5\) (\(A\to BaB\)), но его нельзя применить к \(S\). Значит, правый вывод
    по этой последовательности тоже невозможен.

    \item Эквивалентные левый и правый выводы с индексами правил

    Ту же терминальную цепочку
    \[
    w=bbbbababbaba
    \]
    можно получить левым и правым выводами, но с другим порядком применения
    правил.
    \begin{enumerate}
 
    \item Левый вывод
    \[
    \begin{aligned}
    S_2 &\Rightarrow bS_2 \Rightarrow bbS_2 \Rightarrow bbbS_1 \Rightarrow bbbA_4ba \\
        &\Rightarrow bbbA_5aSba \Rightarrow bbbB_7aBaSba  \\
        &\Rightarrow bbbbaB_7aSba \Rightarrow bbbbaaS_2ba \\
        &\Rightarrow bbbbaabS_1ba \Rightarrow bbbbaabA_6baba \\
        &\Rightarrow bbbbaababbaba
    \end{aligned}
    \]
    
    Последовательность правил левого вывода:
    \[
    2,\;2,\;2,\;1,\;4,\;5,\;7,\;7,\;2,\;1,\;6.
    \]
    
    \item Правый вывод
    \[
    \begin{aligned}
    S_2 &\Rightarrow bS_2 \Rightarrow bbS_2 \Rightarrow bbbS_1 \Rightarrow bbbA_4ba \\
        &\Rightarrow bbbAaS_2ba \Rightarrow bbbAabS_1ba \Rightarrow bbbAabA_6baba \\
        &\Rightarrow bbbA_5abbaba \Rightarrow bbbBaB_7abbaba \\
        &\Rightarrow bbbB_7ababbaba \Rightarrow bbbbababbaba
    \end{aligned}
    \]
    Последовательность правил правого вывода:
    \[
    2,\;2,\;2,\;1,\;4,\;2,\;1,\;6,\;5,\;7,\;7.
    \]
    \end{enumerate}
    
    \end{enumerate}
